\documentclass[../main.tex]{subfiles}
\begin{document}
%===================TIÊU ĐỀ TẬP=============================
\begin{table}[h]
  \begin{adjustwidth}{-1cm}{}
    \begin{tabular}{>{\centering\arraybackslash}p{6cm}|>{\raggedright\arraybackslash}p{15cm}}
      \multirow{5}{6cm}
      % Cột bên trái: ÉPISODE và số tập
      {\\[-40pt]\flaregothic\fontsize{20pt}{20pt}\selectfont \centering
        \color{eptype}ÉPISODE\color{black}\\
        \burbank\fontsize{72pt}{72pt}\selectfont
        \color{epnum}-10\color{black}\\[6pt]
        \cafeteria\fontsize{20pt}{20pt}\selectfont\color{epnum}(P-1)\color{black}\\
      }
       &
      % Dòng thứ nhất: tên tập tiếng Nhật
      {\fotskip\fontsize{18pt}{18pt}\selectfont
        おれさまのごしゅじん}
      \\[6pt]
      % Dòng thứ hai: tên tập tiếng Anh
       & {\cafeteria\fontsize{18pt}{18pt}\selectfont My Boss}                           \\[4pt]
      % Dòng thứ ba: tên tập tiếng Việt
       & {\vnmsans\fontsize{18pt}{18pt}\selectfont Một ngày của Sakura Miko}            \\[8pt]
      % Dòng thứ tư: Nhân vật xuất hiện
       & {\caviar{\fontsize{14pt}{14pt}\selectfont Nhân vật: \emp{miko} \emp{kintoki}}} \\[8pt]
      % Dòng cuối cùng: Thời gian diễn ra của tập (thường sẽ là ngày phát hành)
       & {\continuum\fontsize{16pt}{16pt}\selectfont Ngày phát hành: 07/12/2018}
    \end{tabular}
  \end{adjustwidth}
\end{table}
%========================VĂN BẢN==============================
\color{black}

\txtbx[2]{
  {\color{vol} \uline{Tóm tắt:}} Mở đầu cho quyển tiền truyện là những lời tự sự đầy bất lực của chú mèo thần Kintoki về cô chủ Sakura Miko của mình. Bất chấp việc sở hữu sức mạnh phép thuật thực sự, Miko lại thích dành cả ngày để làm những trò ngốc nghếch, vô tư lự, khiến vị thức thần cao quý phải ngậm ngùi chấp nhận số phận làm "bảo mẫu" cho một cô nàng Vu nữ tự xưng là "ưu tú" nhưng lại lười biếng khó ai sánh bằng.
}

\pov{Kintoki}

Tên của ta là Kintoki. Ta là một thức thần mang hình hài của một chú mèo hồng biết bay, sở hữu đôi mắt vàng sắc lẹm cùng một chiếc đai cổ đeo chuông lục lạc vàng bóng loáng. Ta tự hào là một linh thú cao quý, mang trong mình trọng trách bảo vệ sự linh thiêng và thanh tịnh cho chốn đền thờ này.

\img{p1-1a}

Và kia, cái người mà ta luôn phải ngậm đắng nuốt cay gọi một tiếng ``cô chủ'', bằng một cách thần kỳ nào đó đã thành công triệu hồi ta đến với Ngôi đền Ảo này. Đúng rồi đấy, chính là cái cô ả tóc hồng đang loi nhoi vắt vẻo trên lan can cây cầu gỗ đỏ đằng kia kìa.

Ban đầu, khi mới được triệu hồi tới đây, ta đã ôm một niềm huyễn hoặc lớn lao rằng bản thân sẽ được kề vai sát cánh cùng một bậc vu nữ đúng nghĩa: một người đĩnh đạc, uy nghiêm, nói năng cẩn trọng và toát ra thứ khí chất tâm linh thanh cao. Nhưng thực tế thì sao? Nhìn cô ta đi, các ngươi thấy có giống một vu nữ không?

\img{p1-1b}

Thôi được rồi, để ta tóm tắt lại tình cảnh thảm thương của mình cho các ngươi nghe. Đây là Ngôi đền Ảo (\textit{Virtual Shrine}), nơi ta vốn dĩ cai quản và quán xuyến mọi thứ, chí ít là trước khi bị con nhóc vu nữ ngốc nghếch kia lôi cổ tới đây. Từ một Thức thần quyền uy chuyên trừ tà diệt ma, giờ đây ta lại bị giáng cấp xuống làm một gã ``bảo mẫu'' bất đắc dĩ cho một cô nàng lúc nào cũng tự vỗ ngực xưng tên là ``Vu nữ Ưu tú''. Thật đúng là một trò đùa thế kỷ của tạo hóa.

Nhìn cái điệu bộ ngồi vắt vẻo trên thành lan can với gương mặt ngông nghênh tự mãn của cô ta kìa! Cái sự tự tin thừa thãi đó là từ đâu ra vậy? Cứ làm như thể bản thân là trung tâm của cả vũ trụ vậy!

Thế nhưng, có một sự thật thú vị mà ta buộc phải thừa nhận: cô ta \textit{thực sự} có sức mạnh tâm linh. Và đó không phải là thứ ``sức mạnh tinh thần thanh cao'' mang tính trừu tượng mà đám pháp sư hay lòe bịp thiên hạ đâu. Sức mạnh của cô ta là thật, theo đúng nghĩa đen của nó.

\begin{figure}[H]
  \centering
  \begin{subfigure}[t]{0.45\textwidth}
    \centering
    \includegraphics[width=8cm]{../../../../Image/BookP/p1-1c1.png}
  \end{subfigure}
  \quad
  \begin{subfigure}[t]{0.45\textwidth}
    \centering
    \includegraphics[width=8cm]{../../../../Image/BookP/p1-1e.png}
  \end{subfigure}
\end{figure}

Những tia sét ma thuật màu xanh biếc kia hoàn toàn không phải là ảo ảnh; chúng liên tục chạy dọc quanh cơ thể cô ta một cách cực kỳ tự nhiên. Hơn thế nữa, cô ta còn có thể tùy ý phóng ra những luồng năng lượng hình lục giác chói lòa ấy khỏi tay mình. Nghe ngầu lắm đúng không? Nhưng các ngươi có biết cô ta dùng luồng sức mạnh hủy diệt đó để làm gì không? Chẳng để làm cái quái gì cả!

Nhớ có lần, ta tò mò hỏi: ``Miko, cô định dùng nguồn năng lượng đó vào việc gì vậy?'', và đoán xem cô ả đã trả lời ta bằng cái giọng điệu hồn nhiên đến mức nào:  ``Ể? Thì để trang trí thôi mà, trông lấp lánh đẹp mắt mà, đúng không ne!?''

Nghe có tức lộn ruột không cơ chứ!? Nhưng biết làm sao được, đó chính là Sakura Miko - cô nàng vu nữ kỳ quái và phí phạm tài năng nhất mà ta từng gặp trong suốt cuộc đời làm mèo thần của mình.

\ct

Dạo gần đây, tần suất cô ta hành xử quái lạ ngày càng gia tăng đến mức báo động. Chẳng hạn như việc cô ta ngang nhiên vác cây gậy \textit{gohei} đi trừ tà chạy nhông nhông khắp khuôn viên đền, không ngừng vung vẩy múa may quay cuồng rải những luồng sét xanh bay tứ tung, vừa làm vừa cười toe toét một cách vô cùng mãn nguyện. Các ngươi thử đặt mình vào vị trí của ta xem, nhìn cảnh đó mà không lộn ruột thì mới là lạ đấy! Vu nữ kiểu gì mà cư xử chẳng khác nào một đứa trẻ con lên ba đang hớn hở chơi đồ hàng vậy chứ?

\img{p1-1c2}

À, ta còn nghe ngóng được rằng dạo này cô ta vừa mới dấn thân vào một lĩnh vực hoàn toàn mới - làm một \textit{Virtual YouTuber}. Đúng vậy đấy, cô vu nữ tự xưng này bắt đầu mày mò lên mạng để phát sóng trực tiếp (\textit{live-stream}) chơi game. Chắc các ngươi đang tò mò tự hỏi: ``Thế số lượng người xem có đông đảo không?'' Hừm\dots\ lẹt đẹt và ế ẩm thảm thương, tới mức số người bình luận đôi khi còn đếm trên đầu ngón tay.

Lẽ ra, theo tư duy logic thông thường, cô ta phải biết dồn toàn tâm toàn ý lo cho công việc nhang khói và tế lễ ở đền chùa trước đã. Nhưng không, ``cô chủ'' yêu quý của ta lại chọn cách ``ăn chơi đua đòi'' chạy theo xu hướng trên mạng ảo như thế đấy.

Cô ả vừa híp mắt cười tủm tỉm, vừa hớn hở giơ miếng bánh cá hình con taiyaki trước mặt ta, giọng nũng nịu mời mọc: ``Kintoki! Cho ông ít bánh cá này! Há mồm ra coi\~{}''

\img{p1-1d}

Chê bai thì chê bai thế thôi, nhưng thỉnh thoảng cô ta cũng ra dáng một người chủ tử tế. Cứ mỗi lần thấy cô ả hồ hởi giơ miếng bánh cá hình con taiyaki trước mặt, ta biết ngay giờ ăn đã đến. Dẫu vậy, ta chẳng dại gì mà ngoạm ngay lấy miếng cá trên tay cô ả. Ai biết được với cái tính khí thất thường kia, cô ta có đang định giở trò đùa dai gì không chứ? Cẩn tắc vô áy náy!

``Quả nhiên làm vu nữ đúng là biểu tượng của sự ưu tú thật nha,'' cô ả khẽ thở dài thườn thượt, khuôn mặt lộ rõ vẻ uể oải giả trân như thể vừa lao lực ghê gớm lắm. ``Nhưng mà làm đúng là mệt thật đấy. Ông cũng nghĩ thế phải hơm, Kintoki?''

Trời đất thánh thần ơi, nghe đến đây, ta chỉ muốn bay lại cào thẳng vào mặt cô ta một trận: ``Mới làm được ba cái trò mèo vớ vẩn mà đã dám mở miệng kêu mệt!?'' Nhưng, nghĩ đi nghĩ lại, ta đành ngậm bồ hòn làm ngọt mà quyết định im lặng. Dù sao thì ở cái chốn đền thờ này, ta cũng được xem như một ``người cha đỡ đầu'' già dặn kinh nghiệm cơ mà. Nhưng thú thật là nó phiền phức lắm, đã vụng về vô dụng lại còn kênh kiệu ảo tưởng, lúc nào cũng cư xử y như thể mình là cái rốn của vũ trụ vậy!

\ct

Một trong những trò tiêu khiển thường ngày khác của Sakura Miko, ngoài việc vung vẩy gậy trừ tà, là hùa cùng với ông Thần rảnh rỗi sinh nông nổi (\textit{Kami-sama}) để chế cháo ra đủ thứ phép thuật kỳ quặc. Ta thề là ta hoàn toàn không biết mục đích thực sự của họ là gì, nhưng ta dám cá cược bằng cả cái chuông vàng này rằng nó chắc chắn chẳng phải việc gì đứng đắn hay nghiêm túc đâu.

À mà này, ta đã kể cho các ngươi nghe chưa nhỉ? Cô ả mít ướt lắm. Mong manh dễ vỡ cực kỳ. Chỉ cần một chút chuyện nhỏ nhặt không diễn ra đúng như ý muốn thôi là cái môi mỏng dính kia đã trề ra, mặt mày nhăn nhó như sắp khóc tu tu đến nơi rồi.

Sau một hồi tự xưng là bản thân đã ``làm việc chăm chỉ đến mức kiệt sức'' - đúng, ta đang dùng một tông giọng đậm mùi mỉa mai đấy - ``cô chủ'' lại quay sang nhìn ta với một vẻ mặt mãn nguyện đến khó tả: ``Không ổn rồi, tôi mệt mỏi rã rời luôn á\dots\ Kintoki, ta cùng vào trong phòng nghỉ ăn vặt tí nha\~{}''

Nói dứt lời, cô ả nở một nụ cười mà theo ta nghĩ là rất ``đậm chất thảo mai''. Nhìn cái điệu bộ hớn hở đó, ai không biết lại tưởng cô ta đang chờ đợi những lời tung hô tán thưởng từ ta cơ đấy. Thật đúng là làm như thể bản thân vừa hoàn thành một chiến dịch giải cứu thế giới vĩ đại vậy!

Thật tình, ta tự hỏi, cái danh xưng \textit{Elite} (Ưu tú) này rốt cuộc là từ trên trời rơi xuống hay sao? Cô ả quay xong video tự sướng về cuộc sống của mình chưa vậy? Liệu cái đầu nhỏ bé màu hồng kia có thực sự hiểu được cái định nghĩa về ``công việc của một vu nữ'' không?

Dường như bắt bài được ánh mắt chất vấn đầy khinh bỉ của ta, cô ả lập tức đanh giọng lại lấp liếm: ``Im đi nha! Ông phải biết tôi cũng làm việc vất vả và xứng đáng được nghỉ ngơi chứ! Cẩn thận kẻo tôi cắt quà bánh của bây giờ!''

Ngon, dọa dẫm luôn cơ đấy. Làm YouTube nhiều chắc Sakura Miko kia cũng tiêu hao nhiều calo lắm rồi nhỉ? Ta chỉ biết thở dài lắc đầu, chẳng buồn mở miệng đôi co cãi lại. Thôi thì, cứ rộng lượng chiều lòng ``cô chủ'' một chút vậy. Cứ để cô ả tận hưởng niềm vui nhỏ nhoi với đống bánh cá của mình đi.

Nói xấu nhiều như thế, bóc phốt mỏi cả mồm như thế, nhưng nếu bảo rằng ta không hề gắn bó với cô ta thì quả là một lời nói dối trắng trợn. Thú thật, nếu không có cái sự ồn ào của Sakura Miko, ta cũng chẳng biết phần đời còn lại của mình sẽ trôi dạt về đâu trong cái Ngôi đền Ảo hiu quạnh này. Đôi khi ta không tài nào hiểu nổi lối tư duy kỳ quặc và lối sống bất cần đời của cô ta. Cô ta cứ lơ ngơ, hành động bốc đồng theo ý thích mà chẳng thèm quan tâm đến ánh mắt của người đời. Không biết sau này, với cái tính cách ngốc nghếch đó, cô ả có làm nên cơm cháo gì trong cái thế giới ngoài kia không đây\dots

Nhưng, nghĩ đi nghĩ lại, mặc cho việc cô ta cứ thích làm màu, thích vung vẩy cái gậy \textit{gohei} tạo phép khắp nơi như thể biến nơi đây thành sân khấu riêng của mình, ta đã quen với cái sự ồn ào đáng yêu đó rồi. Có lẽ, chính cái sự ngốc nghếch, vô tư và phiền phức mang tên Sakura Miko ấy mới là ngọn lửa sưởi ấm, giữ cho ngôi đền này luôn tràn ngập tiếng cười và sự sống động.

\end{document}